%!TEX TS-program = xelatex
\DocumentMetadata{
  pdfversion=2.0,
  lang=es-MX,
  pdfstandard=ua-2
}

\documentclass{sener2025}

\title{Prueba de notas al pie en \textbackslash fuente}
\author{SENER}
\date{\today}

\begin{document}
\portadafondo

\section{Prueba de Reproducción Exacta}

Para el año 2024, el consumo no energético representó el 3.3\% (183.18 PJ) del consumo final en el país. Predominando otros energéticos con 52.3\% (95.74 PJ), seguidos por gas seco con 41.62\% (76.25 PJ), gasolinas y naftas con 5.83\% (10.69 PJ). De forma complementaria, se utilizan GLP y bagazo de caña, como se observa en la Figura 5.19.

La relevancia de otros energéticos en el consumo final no energético se ha mantenido durante el periodo 2010-2024 (Figura 5.20.). Sin embargo, el gas seco ha registrado un crecimiento sostenido durante los últimos 5 años. De hecho, de 2023 a 2024, mientras el consumo de otros energéticos cayó 9\%, el gas seco creció 35\%.

\begin{figure}[H]
  {\centering
  % Texto alternativo para accesibilidad
  \pdftooltip{\rule{10cm}{5cm}}{Figura 5.19 Composición del consumo final no energético por energético, 2024}
  \par}
  \raggedright
  \caption{Figura 5.19 Composición del consumo final no energético por energético, 2024}
  \label{fig:figura_5_19_composicion_del_co}
\end{figure}
\vspace{-4pt}
\fuente{Elaboración SENER. \footnote{Otros energéticos contemplan a lubricantes, parafinas, asfaltos, azufre, materia prima negro de humo, gasóleo, etano, gas de alto horno, gas de coque, licor negro, gas residual y aceite residual.}}

\end{document}
